% !TeX spellcheck = en_US
\chapter{State of the Art}\label{ch:sota}%

This chapter has two parts. Section~\ref{sec:theory} introduces the theoretical concepts the thesis builds on: reinforcement learning with privileged critics, imitation learning and its compounding-error problem, interactive data aggregation, generative action heads based on diffusion and flow matching, visual representations for manipulation, and operational-space control. Section~\ref{sec:related} reviews related work on robotic assembly with dexterous hands, policy chaining for long-horizon tasks, sim-to-real transfer, and visuomotor policy distillation.

\section{Theoretical Background}\label{sec:theory}

\subsection{Reinforcement Learning with Privileged Information}\label{sec:rl-background}

Reinforcement learning formalizes sequential decision making as a Markov decision process (MDP) $(\mathcal{X}, \mathcal{A}, P, r, \gamma)$ with state space $\mathcal{X}$, action space $\mathcal{A}$, transition kernel $P(\vec{x}_{t+1}\mid \vec{x}_t,\vec{a}_t)$, reward $r(\vec{x}_t,\vec{a}_t)$, and discount factor $\gamma\in[0,1)$~\cite{Kober2013}. A policy $\pi(\vec{a}\mid\vec{x})$ is sought that maximizes the expected discounted return $J(\pi)=\E_\pi[\sum_t \gamma^t r_t]$. For robot manipulation the full state is not observable at deployment; the problem is then a partially observable MDP (POMDP), in which the policy acts on observations $\vec{z}_t\sim O(\cdot\mid\vec{x}_t)$ instead of states.

Proximal policy optimization (PPO)~\cite{Schulman2017} is the on-policy actor--critic algorithm used throughout this thesis. The actor $\pi_\theta$ is updated by maximizing the clipped surrogate objective
\begin{equation}\label{eq:ppo}
  L^{\mathrm{PPO}}(\theta) = \E_t\!\left[\min\!\left(\rho_t(\theta)\hat{A}_t,\; \operatorname{clip}\!\left(\rho_t(\theta),1-\epsilon,1+\epsilon\right)\hat{A}_t\right)\right],
  \qquad \rho_t(\theta)=\frac{\pi_\theta(\vec{a}_t\mid\vec{z}_t)}{\pi_{\theta_{\mathrm{old}}}(\vec{a}_t\mid\vec{z}_t)},
\end{equation}
where $\hat{A}_t$ is an advantage estimate computed from the critic's value prediction. The clipping bounds the policy change per update and makes PPO stable on the thousands of parallel environments that GPU simulators provide.

In an asymmetric actor--critic the actor and the critic receive different inputs~\cite{Pinto2017}. The critic estimates the value function from the complete simulator state, including quantities that no sensor measures, whereas the actor observes only a noisy subset that can later be reproduced on hardware or distilled into a sensor-based policy. Because the critic is discarded after training, the privileged information shapes the gradients without entering the deployed policy. When the actor observation is itself a noisy or delayed function of the state, a recurrent actor with a long short-term memory (LSTM) cell~\cite{Hochreiter1997} integrates the observation history into a hidden state $\vec{h}_t$ and thereby recovers part of the information lost to partial observability. This recipe, together with domain randomization of physical and controller parameters, underlies most recent sim-to-real results in dexterous manipulation~\cite{Handa2024,Petrenko2023,Lum2024,Singh2025}.

\subsection{Imitation Learning and Compounding Error}\label{sec:il-background}

Imitation learning (IL) trains a policy from demonstrations rather than from a reward. Given a dataset $\Dset=\{(\obs_t,\act^{*}_t)\}$ of observations and expert actions, behavior cloning (BC) fits a policy by supervised regression,
\begin{equation}\label{eq:bc}
  \theta^{*}=\argmin_\theta \E_{(\obs,\act^{*})\sim\Dset}\!\left[\ell\!\left(\piS(\obs),\act^{*}\right)\right],
\end{equation}
with a per-sample loss $\ell$ such as the squared error. BC is simple and data-efficient, but it is trained on the state distribution induced by the expert and evaluated on the state distribution induced by the student. Small prediction errors move the student into states the expert never visited, where its predictions are unreliable, and the error grows with the horizon. Ross et al.~\cite{Ross2011} show that the expected cost of a BC policy with per-step error $\varepsilon$ can grow as $O(T^2\varepsilon)$ over a horizon $T$, because an error at step $t$ affects all subsequent steps. For a cyclic task of several hundred control steps this quadratic growth is the central obstacle: a policy that reproduces each phase well in isolation can still fail after a few cycles once small deviations in grasp height or nut angle accumulate.

Dataset aggregation (DAgger)~\cite{Ross2011} addresses the distribution mismatch by querying the expert on the states that the student itself visits. In iteration $i$ the environment is driven by a mixture policy that executes the expert's action with probability $\beta_i$ and the student's action otherwise, the expert's action is recorded as the label for every visited state, the labeled states are added to the aggregated dataset, and the student is retrained. Under a no-regret learner, the resulting policy has expected cost $O(T\varepsilon)$, linear rather than quadratic in the horizon. The guarantee holds because the aggregated dataset converges to the state distribution of the learned policy itself, so that the supervised error measured on that dataset is the error the policy will actually incur; behavior cloning measures its error on the expert's distribution, where the policy is never evaluated. The price is the $O(T)$ expert queries per iteration and the requirement that the expert be queryable at every visited state. The schedule $\beta_i$ is typically decayed from one toward zero so that early iterations stay close to the expert's distribution and later iterations expose the student to its own mistakes. SafeDAgger~\cite{Zhang2016} reduces the number of expert queries by invoking the expert only when a learned safety classifier, or a threshold on the disagreement between expert and student actions, predicts that the student is about to fail. The predecessor of this thesis applied DAgger and SafeDAgger to distill a privileged grasping teacher into an RGB student and found that the safety-gated variant produced both the best grasp success and the fewest unsafe episodes~\cite{Oertel2026}.

DAgger requires an expert that can be queried at arbitrary student-visited states. A privileged simulation policy satisfies this requirement, a recorded demonstration does not. The hybrid teacher of this thesis is such a queryable expert, and the method of Chapter~\ref{ch:method} uses this property in two ways: offline, by disturbing the teacher during collection so that it labels its own recoveries from off-nominal states, and online, through a DAgger variant with a flow-matching head that serves as a comparison.

\subsection{Generative Action Heads: Diffusion and Flow Matching}\label{sec:generative-background}

A deterministic regression head predicts one action for each observation. Demonstrations of manipulation tasks are, however, often multimodal: two demonstrators, or one teacher recovering from two different disturbances, may approach the same nut from different sides, and averaging their actions produces an action that neither of them would take. Generative action heads represent the full conditional distribution $p(\act\mid\obs)$ and sample from it at run time. Diffusion Policy~\cite{Chi2024} models the distribution of action sequences by iterative denoising: Gaussian noise is added to demonstrated action chunks in a forward process, a network is trained to predict the noise, and actions are generated by reversing the process over tens to hundreds of denoising steps. Diffusion policies handle multimodality and predict action chunks, which improves temporal consistency, but their sampling cost is high for control loops with tight latency budgets.

Flow matching~\cite{Lipman2023,Lipman2024} reaches the same goal with a simpler training target and a cheaper sampler. Instead of a stochastic denoising process it learns a time-dependent velocity field $\vec{v}_\theta(\vec{x},\lambda)$ that transports samples of a simple source distribution $p_0$, typically a standard Gaussian, to the data distribution $p_1$ along the ordinary differential equation (ODE) $\dot{\vec{x}}=\vec{v}_\theta(\vec{x},\lambda)$ for $\lambda\in[0,1]$. The conditional flow-matching objective regresses the velocity field onto the velocity of a prescribed conditional probability path between a source sample and a data sample. With the linear interpolant $\vec{x}^{\lambda}=(1-\lambda)\vec{x}^{0}+\lambda\vec{x}^{1}$, whose velocity is $\vec{x}^{1}-\vec{x}^{0}$, the objective is a plain regression and the learned ODE can be integrated with a handful of Euler steps. Section~\ref{sec:flow-matching} derives this objective and states it for conditional action chunks. Diffusion and flow matching are closely related. The forward noising process of a diffusion model defines a Gaussian probability path between data and noise, and its reverse process can equally be written as a deterministic probability-flow ODE driven by the learned score; flow matching with Gaussian conditional paths contains this diffusion path as a special case, and with the optimal-transport path it replaces the curved diffusion trajectories by straight lines that are cheaper to integrate~\cite{Lipman2023}. The practical consequences are the absence of a noise schedule, a regression target that is a simple difference of two samples, and samplers with few steps. Flow matching has become the action head of choice for large vision-language-action models such as $\pi_0$~\cite{Black2026}, which generate 50-step action chunks with ten integration steps, and it is the action head of the student in this thesis.

\subsection{Visual Representations for Manipulation}\label{sec:vision-background}

A visuomotor student must extract the nut pose relative to the hand from camera data. Three families of encoders are relevant here. Image-space encoders apply a convolutional backbone such as ResNet-18~\cite{He2016} to the RGB image, possibly with depth as an additional channel, and pool the feature map into a vector. They inherit ImageNet pretraining but have to learn metric geometry from appearance and must be made robust to texture and lighting with visual domain randomization~\cite{Tobin2017}. Point-cloud encoders back-project the depth image through the calibrated camera model into a set of 3D points, optionally attach the pixel color, and encode the set with a small shared point network followed by global max pooling in the spirit of PointNet~\cite{Qi2017}. 3D Diffusion Policy (DP3)~\cite{Ze2024} shows that such a compact point encoder generalizes better across viewpoints and appearance than image encoders for diffusion-based visuomotor policies, because the metric geometry enters the representation directly. Adapt3R~\cite{Wilcox2025} combines the two: dense 2D features from a frozen pretrained backbone are localized in 3D by the depth map and expressed in the end-effector frame before attention pooling, which makes the representation invariant to the embodiment and robust to domain change. All three families are implemented in the code base of this thesis, and Section~\ref{sec:student} explains why the DP3-style encoder was retained.

\subsection{Operational-Space Control}\label{sec:osc-background}

Both the teacher and the student command a Cartesian end-effector target that a task-space impedance controller converts to joint torques. Let $\vec{q}\in\R^7$ be the arm joint positions, $\mat{J}(\vec{q})$ the geometric Jacobian of the controlled frame, and $\vec{e}\in\R^6$ the position and axis-angle orientation error between target and current pose. The controller applies
\begin{equation}\label{eq:osc}
  \vec{\tau}=\mat{J}^{\mathrm{T}}\!\left(\mat{K}_p\vec{e}-\mat{K}_d\dot{\vec{x}}\right)+\mat{N}^{\mathrm{T}}\!\left(\mat{K}_{p,0}(\vec{q}_0-\vec{q})-\mat{K}_{d,0}\dot{\vec{q}}\right),
\end{equation}
where $\dot{\vec{x}}$ is the end-effector twist, $\mat{K}_p$ and $\mat{K}_d=2\sqrt{\mat{K}_p}$ are diagonal stiffness and critical-damping gains, and the second term pulls the redundant seventh degree of freedom toward a rest posture $\vec{q}_0$ through a nullspace projector $\mat{N}$. The projector is dynamically consistent when $\mat{N}^{\mathrm{T}}=\mat{I}-\mat{J}^{\mathrm{T}}\bar{\mat{J}}^{\mathrm{T}}$ with $\bar{\mat{J}}=\mat{M}^{-1}\mat{J}^{\mathrm{T}}(\mat{J}\mat{M}^{-1}\mat{J}^{\mathrm{T}})^{-1}$ and joint-space inertia $\mat{M}$~\cite{Dietrich2015,RoboticSystemsLab2019}. Equation~\eqref{eq:osc} is the Factory and FORGE controller~\cite{Narang2022,Noseworthy2025}, which this thesis inherits. Premultiplying the task-space force by the task-space inertia $\mat{\Lambda}=(\mat{J}\mat{M}^{-1}\mat{J}^{\mathrm{T}})^{-1}$ yields the full operational-space formulation, in which the closed loop becomes a decoupled unit-mass system in every direction. Impedance control is preferred over stiff position control for a contact-rich task because the compliant law yields on unexpected contact rather than driving through it, and because the policy can probe the impedance by commanding pose offsets rather than torques. The compliance is also what allows the simplified thread model of Section~\ref{sec:thread-model} to tolerate small deviations of the hand from the helical path.

\section{Related Work}\label{sec:related}

\subsection{Robot Assembly with Dexterous Manipulation}\label{sec:related-assembly}

Assembly tasks such as peg-in-hole insertion, gear meshing, connector mating, screwdriving, and nut--bolt fastening demand precise motion and closed-loop contact under geometric and dynamic uncertainty~\cite{Siciliano2016}. Industrial practice solves them with fixtures and dedicated tools. A multi-fingered hand removes the need for task-specific tooling but brings high-dimensional control and complex contact dynamics with it~\cite{BicchiKumar2000,Weinberg2024}. On the simulation side, Factory~\cite{Narang2022} made nut-and-bolt assembly learnable by combining signed-distance-field collision detection, contact reduction, and a Gauss--Seidel solver, so that a thousand parallel assemblies run in real time on one GPU; it also splits the task into pick, place, and screw subpolicies. FORGE~\cite{Noseworthy2025} adds force-guided contact search under pose uncertainty, and IndustReal~\cite{Tang2023industreal} transfers simulation-trained assembly policies to a physical arm. All three use a parallel-jaw gripper. For dexterous hands, SimToolReal~\cite{Kedia2026} learns zero-shot tool use from simulation, and Play2Perfect~\cite{Lum2026} pretrains a goal-conditioned manipulation policy through play and fine-tunes it for precise assembly, including screwing a rod-like furniture leg; this gives goal-directed screwing in a structured setting, but not for everyday threaded parts such as nuts. Screwdriving with hand tools has been modeled analytically and learned with physics-informed models~\cite{Tang2023,Tang2024screwdriving,Manyar2023}, and finger gaiting under contact-force control has been shown for pivoting~\cite{Xue2023}. The closest prior work is DexScrew~\cite{Hsieh2026}, which uses a hand--arm setup comparable to the one here and replaces the thread geometry by a revolute-joint proxy to learn a rotational finger gait; the gait assists teleoperation and supports the collection of real demonstrations for tactile behavior cloning. Its policy spins the fastener without a target angle and without learned termination. Goal-directed threading over many grasp-and-release cycles, which this thesis addresses for several nut sizes, is left open by these works.

\subsection{Policy Chaining in Long-Horizon Tasks}\label{sec:related-chaining}

A single policy for a long-horizon task faces two difficulties. Small control errors compound over the horizon~\cite{Ross2011}, and fine-tuning the policy on later subtasks can erase what it learned on earlier ones, the catastrophic forgetting problem~\cite{Kirkpatrick2017}. Larger networks, frozen or regularized parameters, and architectures with longer temporal memory such as recurrent networks and Transformers mitigate these effects but do not remove them. The alternative is to decompose the task into skills with their own policies and to chain them. Each skill is then easier to learn, but the chain introduces a transition problem: the states in which one skill terminates may lie outside the initial-state distribution on which the next skill was trained, so the chain breaks at the boundary. Adversarial skill chaining~\cite{Lee2022} regularizes the terminal-state distribution of each skill toward the initial-state distribution of its successor. Sequential Dexterity~\cite{Chen2023} learns a transition-feasibility function with which adjacent dexterous policies are fine-tuned, switches between them autonomously, and recovers from failed transitions. The method of this thesis decomposes only during data generation, where every phase has a simple specification, and unifies at deployment by distilling the chained demonstrations into one policy; the transitions are then part of what the student learns rather than a problem to be solved between separate policies.

\subsection{Sim-to-Real Transfer}\label{sec:related-sim2real}

Policies trained in simulation meet discrepancies in perception, dynamics, contact, and control when they are deployed on hardware, and the discrepancies grow with the dimensionality of the control problem and the amount of contact involved. The most common response is to make the policy robust enough for zero-shot transfer, by randomizing visual and physical parameters during training~\cite{Tobin2017,Handa2024} or by making the visual representation itself transferable~\cite{Wilcox2025}. A second family adapts a simulation-trained policy with a small amount of real-world data, for instance by learning a residual policy from on-robot corrections~\cite{Jiang2025transic} or by continuing training on the robot~\cite{Luo2025,Huang2025}. This thesis follows the first family. Its student is trained on simulated data only, and transfer rests on three properties of the training setup: the camera model is calibrated against the physical cell and reproduced in simulation, the visual representation is a metric point cloud in the spirit of DP3~\cite{Ze2024}, which is invariant to appearance up to depth noise, and the rendered appearance is randomized in the style of DextrAH-RGB~\cite{Singh2025}. The approach requires no real-world data collection and no training on the robot.

\subsection{Visuomotor Policy Distillation}\label{sec:related-distillation}

Learning by Cheating~\cite{Chen2020} established the route from a policy that is trained with privileged simulator state to a deployable policy that acts from onboard sensing. In dexterous manipulation, DextrAH-G~\cite{Lum2024} and DextrAH-RGB~\cite{Singh2025} distill a privileged RL grasping teacher into depth- and RGB-based students with online DAgger, and UniGraspTransformer~\cite{Wang2025} consolidates many object-specific teachers into one student. The author's semester thesis~\cite{Oertel2026} applied this recipe to a 23-DoF hand--arm system and compared behavior cloning, DAgger, and SafeDAgger. For long-horizon behavior, two further ingredients matter. Predicting a sequence of future actions instead of a single action improves temporal consistency: Action Chunking with Transformers (ACT)~\cite{Zhao2023} predicts chunks with a Transformer and blends overlapping chunks by temporal ensembling. Representing the action distribution generatively lets a policy keep several valid behaviors apart instead of averaging them: Diffusion Policy~\cite{Chi2024} models multimodal action sequences by iterative denoising, DP3~\cite{Ze2024} conditions it on compact point-cloud features, and flow matching~\cite{Lipman2023} learns a continuous vector field that generates actions by integrating an ODE in few steps, as the vision-language-action model $\pi_0$~\cite{Black2026} does. The student of this thesis combines all three ingredients, a point-cloud encoder, action chunking with temporal ensembling, and a flow-matching head, because its training data mix nominal teacher actions with the teacher's recoveries from disturbances.
